% =====================================================================
%  Banco complementar --- Forca Magnetica  (REVISADO)
%  Substitui a versao gerada por template (5 clones em N1, 10 em N2,
%  10 em N3 e 8 questoes conceituais marcadas como N4).
%  O cap11 ja cobre: condicao de forca nula (MC), condutor de 30 cm com
%  2 A em 0,5 T, regra da mao para carga movendo-se com B entrando no
%  plano, carga elementar a 1e6 m/s em 0,2 T, espira retangular girando
%  (qualitativo), proton em campo uniforme (raio) e o SELETOR DE
%  VELOCIDADES com campos cruzados -- motivo pelo qual é a N4 antiga sobre
%  projeto de seletor foi removida deste banco. Ela e substituida pelo
%  ESPECTROMETRO DE MASSA, que usa o seletor como etapa e vai alem.
%  Este banco explora: forca entre fios e a definicao do ampere,
%  trajetoria helicoidal, torque e momento magnetico, efeito Hall
%  quantitativo, campo nao uniforme, projeto de atuador e a demonstracao
%  de que a forca magnetica nao realiza trabalho.
%  ACRESCIMO POSTERIOR: tres questoes sobre a FORCA DE LORENTZ COMPLETA,
%  F = q(E + v x B). As duas parcelas ja apareciam separadamente e os
%  campos cruzados surgiam no seletor de velocidades e na deriva ExB,
%  mas a expressao unificada nao estava escrita em lugar nenhum.
%  Distribuicao: 5 N1 + 11 N2 + 11 N3 + 9 N4 = 36
% =====================================================================

\subsubsection*{Banco complementar --- Força magnética}

\begin{atencao}[Organização do banco]
Duas expressões governam tudo: $\vec F=q\,\vec v\times\vec B$ para cargas e
$\vec F=I\,\vec\ell\times\vec B$ para condutores. Ambas são \textbf{produtos
vetoriais} --- a força é \emph{perpendicular} tanto à velocidade quanto ao
campo, e se anula quando eles são paralelos. Daí decorre o fato central: a força
magnética \textbf{nunca realiza trabalho}. O N3 trata de trajetórias,
espectrometria e efeito Hall; o N4 exige demonstrações e projeto.
\end{atencao}

\blocofixacao

\begin{questao}[1]{}
Uma carga $q=+\SI{2}{\micro\coulomb}$ move-se a $\SI{500}{\meter\per\second}$
perpendicularmente a um campo de $\SI{0.2}{\tesla}$. Determine a força.
\resposta{$F=\SI{0.2}{\milli\newton}$}
\resolucao{
\[
F=qvB\sin\theta=(\pot{2}{-6})(500)(0{,}2)(1)=\pot{2}{-4}\,\si{\newton}.
\]
\textbf{Direção:} perpendicular \emph{simultaneamente} a $\vec v$ e a
$\vec B$, dada pela regra da mão direita aplicada a
$\vec v\times\vec B$ (e invertida, se a carga for negativa).}
\end{questao}

\begin{questao}[1]{}
Um condutor retilíneo de $\SI{40}{\centi\meter}$ conduz $\SI{3}{\ampere}$
perpendicularmente a $B=\SI{0.25}{\tesla}$. Determine a força.
\resposta{$F=\SI{0.3}{\newton}$}
\resolucao{
\[
F=BIL\sin\theta=(0{,}25)(3)(0{,}40)(1)=\SI{0.3}{\newton}.
\]
\textbf{Origem da fórmula:} ela é a soma das forças sobre todos os portadores do
condutor. Escrevendo $I=nqv_dA$ e $L$ o comprimento, o produto
$nqv_dA\cdot L$ reproduz $qv_d$ somado sobre $nAL$ portadores.\\
A força é transmitida ao \emph{fio} pelas colisões dos portadores com a rede
cristalina --- não são os elétrons que "empurram" o fio diretamente.}
\end{questao}

\begin{questao}[1]{}
A mesma carga da primeira questão move-se agora formando $\ang{30}$ com o
campo. Determine a força.
\resposta{$F=\SI{0.1}{\milli\newton}$}
\resolucao{
\[
F=qvB\sin\ang{30}=(\pot{2}{-4})(0{,}5)=\pot{1}{-4}\,\si{\newton}.
\]
\textbf{Metade} do valor perpendicular.\\
\textbf{Casos-limite:} $\theta=\ang{90}$ dá força máxima; $\theta=0$ ou
$\ang{180}$ (velocidade \emph{paralela} ao campo) dá força \textbf{nula} --- a
carga segue em linha reta, sem sentir o campo.}
\end{questao}

\begin{questao}[1]{}
Uma carga positiva move-se para a \textbf{direita} em uma região onde
$\vec B$ aponta para \textbf{cima}. Determine o sentido da força.
\begin{alternativas}[2]
\item para dentro do plano
\item para fora do plano
\item para cima
\item para a direita
\end{alternativas}
\resposta{(a)}
\resolucao{Pela regra da mão direita para $\vec v\times\vec B$: dedos apontando
no sentido de $\vec v$ (direita), curvando-se para $\vec B$ (cima) --- o polegar
aponta para \textbf{dentro} do plano.\\
\textbf{Verificação:} a força deve ser perpendicular a ambos. "Direita" e
"cima" definem o plano da folha, logo a força só pode ser perpendicular a ele
--- o que elimina (c) e (d) imediatamente.\\
\textbf{Se a carga fosse negativa}, a força seria para fora do plano.}
\end{questao}

\begin{questao}[1]{}
A partir de $F=qvB$, mostre que $\SI{1}{\tesla}$ equivale a
$\SI{1}{\newton\second\per\coulomb\per\meter}$.
\resposta{$\si{\tesla}=\dfrac{\si{\newton}}{\si{\coulomb}\cdot
\si{\meter\per\second}}$}
\resolucao{
\[
B=\frac{F}{qv}
\;\Longrightarrow\;
[\si{\tesla}]=\frac{\si{\newton}}{\si{\coulomb}\cdot\si{\meter\per\second}}
=\frac{\si{\newton}\cdot\si{\second}}{\si{\coulomb}\cdot\si{\meter}}
=\frac{\si{\newton}}{\si{\ampere}\cdot\si{\meter}}.
\]
A última forma, $\si{\newton\per\ampere\per\meter}$, decorre diretamente de
$F=BIL$ --- e é a mais prática para conferir contas com condutores.\\
\textbf{O tesla é uma unidade grande:} $\SI{1}{\tesla}$ exerce
$\SI{1}{\newton}$ sobre um metro de fio conduzindo $\SI{1}{\ampere}$.}
\end{questao}

\blocoaplicacao

\begin{questao}[2]{}
Um fio de $\SI{10}{\centi\meter}$ conduz $\SI{2}{\ampere}$ em um campo de
$\SI{0.2}{\tesla}$, formando $\ang{30}$ com ele. Determine a força.
\resposta{$F=\SI{20}{\milli\newton}$}
\resolucao{
\[
F=BIL\sin\ang{30}=(0{,}2)(2)(0{,}10)(0{,}5)=\pot{2}{-2}\,\si{\newton}.
\]
\textbf{Interpretação do seno:} apenas a componente de $\vec\ell$
\emph{perpendicular} a $\vec B$ contribui. Equivalentemente, o comprimento
\emph{efetivo} é $L\sin\theta=\SI{5}{\centi\meter}$.\\
\textbf{Direção da força:} perpendicular ao plano definido pelo fio e pelo
campo --- e \emph{não} perpendicular apenas ao fio.}
\end{questao}

\begin{questao}[2]{}
Um próton entra perpendicularmente em um campo de $\SI{0.5}{\tesla}$ com
$v=\pot{2}{6}\,\si{\meter\per\second}$.
\begin{itensq}
\item Determine o raio da trajetória.
\item Determine o período.
\end{itensq}
\dados{$e=\pot{1.6}{-19}\,\si{\coulomb}$; $m_p=\pot{1.67}{-27}\,\si{\kilogram}$}
\resposta{$r=\SI{4.18}{\centi\meter}$; $T=\SI{131}{\nano\second}$}
\resolucao{(a) A força magnética é centrípeta:
\[
qvB=\frac{mv^{2}}{r}
\;\Longrightarrow\;
r=\frac{mv}{qB}=\frac{(\pot{1.67}{-27})(\pot{2}{6})}{(\pot{1.6}{-19})(0{,}5)}
=\SI{0.0418}{\meter}.
\]
(b)
\[
T=\frac{2\pi r}{v}=\frac{2\pi m}{qB}
=\frac{2\pi(\pot{1.67}{-27})}{(\pot{1.6}{-19})(0{,}5)}
=\pot{1.31}{-7}\,\si{\second}.
\]
\textbf{Fato notável:} o período \textbf{não depende da velocidade} nem do raio.
Partículas rápidas descrevem círculos maiores, mas levam o mesmo tempo para
completá-los.\\
É esse fato que torna possível o \textbf{cíclotron}: a frequência de aceleração
pode ser fixa, embora a partícula ganhe energia a cada volta.}
\end{questao}

\begin{questao}[2]{}
Um elétron e um próton, com a mesma velocidade, entram no mesmo campo
magnético perpendicularmente. Compare seus raios e períodos.
\resposta{$r_p/r_e=1836$; $T_p/T_e=1836$}
\resolucao{Como $r=mv/(qB)$ e $T=2\pi m/(qB)$, ambos são
\emph{proporcionais à massa} (as cargas têm o mesmo módulo):
\[
\frac{r_p}{r_e}=\frac{T_p}{T_e}=\frac{m_p}{m_e}
=\frac{\pot{1.67}{-27}}{\pot{9.11}{-31}}=1836.
\]
\textbf{Sentidos opostos:} tendo cargas de sinais contrários, as duas partículas
curvam-se para lados opostos --- é assim que se identifica o sinal da carga em
câmaras de detecção.\\
\textbf{Consequência experimental:} para conter elétrons e prótons no mesmo
aparelho, o raio do próton é quase dois mil vezes maior. Aceleradores de prótons
são, por isso, muito maiores que os de elétrons para a mesma energia.}
\end{questao}

\begin{questao}[2]{}
Dois fios paralelos longos, separados por $\SI{5}{\centi\meter}$, conduzem
$\SI{10}{\ampere}$ cada. Determine a força por unidade de comprimento.
\dados{$\mu_0=4\pi\times10^{-7}\,\si{\tesla\meter\per\ampere}$}
\resposta{$F/L=\SI{0.4}{\milli\newton\per\meter}$}
\resolucao{Um fio cria campo $B=\dfrac{\mu_0I_1}{2\pi d}$ na posição do outro,
que sofre $F=BI_2L$:
\[
\frac{F}{L}=\frac{\mu_0I_1I_2}{2\pi d}
=\pot{2}{-7}\frac{(10)(10)}{0{,}05}=\pot{4}{-4}\,\si{\newton\per\meter}.
\]
\textbf{Sentido:} correntes de \emph{mesmo} sentido \textbf{atraem-se};
de sentidos opostos, \textbf{repelem-se}.\\
\textbf{Cuidado com a intuição elétrica:} cargas de mesmo sinal se repelem, mas
correntes de mesmo sentido se atraem. As duas situações não são análogas.}
\end{questao}

\begin{questao}[2]{}
Uma espira retangular de área $\SI{100}{\centi\meter\squared}$ com $50$ voltas
conduz $\SI{2}{\ampere}$ em um campo de $\SI{0.3}{\tesla}$.
\begin{itensq}
\item Determine o momento de dipolo magnético.
\item Determine o torque máximo.
\end{itensq}
\resposta{$m=\SI{1}{\ampere\meter\squared}$; $\tau=\SI{0.3}{\newton\meter}$}
\resolucao{(a) $A=\SI{100}{\centi\meter\squared}=\pot{1}{-2}\,\si{\meter\squared}$:
\[
m=NIA=50(2)(\pot{1}{-2})=\SI{1}{\ampere\meter\squared}.
\]
(b)
\[
\tau=mB\sin\theta
\;\Longrightarrow\;
\tau_{\max}=mB=(1)(0{,}3)=\SI{0.3}{\newton\meter}.
\]
\textbf{Quando ocorre o máximo:} com o plano da espira \emph{paralelo} ao campo,
isto é, $\vec m$ perpendicular a $\vec B$. O torque é \textbf{nulo} quando
$\vec m$ está alinhado com $\vec B$ --- posição de equilíbrio.\\
\textbf{Analogia útil:} a espira comporta-se como uma agulha de bússola, e
$\vec m$ faz o papel do "norte" dela.}
\end{questao}

\begin{questao}[2]{}
Explique por que a força magnética não altera o módulo da velocidade de uma
partícula carregada.
\resposta{Ela é sempre perpendicular a $\vec v$, logo não realiza trabalho}
\resolucao{A potência entregue por uma força é $P=\vec F\cdot\vec v$. Como
\[
\vec F=q\,\vec v\times\vec B
\;\Longrightarrow\;
\vec F\perp\vec v
\;\Longrightarrow\;
\vec F\cdot\vec v=0.
\]
Sem potência, não há variação de energia cinética --- e portanto o \emph{módulo}
da velocidade permanece constante.\\
\textbf{O que a força magnética faz:} muda a \emph{direção} do movimento, e só
isso. Ela curva a trajetória sem acelerar nem frear.\\
\textbf{Consequência prática:} um campo magnético sozinho jamais acelera
partículas. Aceleradores usam campos \emph{elétricos} para ganhar energia e
campos magnéticos apenas para \emph{guiar} o feixe.}
\end{questao}

\begin{questao}[2]{}
Um elétron é abandonado \textbf{em repouso} em uma região onde coexistem
$E=\SI{500}{\volt\per\meter}$ e $B=\SI{0.02}{\tesla}$, perpendiculares entre
si.
\begin{itensq}
\item Determine a força e a aceleração no instante inicial.
\item Descreva qualitativamente o movimento subsequente.
\end{itensq}
\dados{$e=\pot{1.6}{-19}\,\si{\coulomb}$; $m_e=\pot{9.11}{-31}\,\si{\kilogram}$}
\resposta{$F=\SI{8e-17}{\newton}$ (só elétrica);
$a=\pot{8.8}{13}\,\si{\meter\per\second\squared}$}
\resolucao{(a) A força total é a \textbf{força de Lorentz}:
\[
\vec F=q\left(\vec E+\vec v\times\vec B\right).
\]
Com $\vec v=\vec 0$, a parcela magnética \textbf{se anula} --- resta apenas a
elétrica:
\[
F=eE=(\pot{1.6}{-19})(500)=\pot{8}{-17}\,\si{\newton};
\qquad
a=\frac{F}{m_e}=\pot{8.8}{13}\,\si{\meter\per\second\squared}.
\]
(b) \textbf{Movimento subsequente.} Assim que o elétron adquire velocidade, a
parcela magnética \emph{deixa de ser nula} e passa a encurvar a trajetória. A
força elétrica continua acelerando; a magnética continua desviando.\\
O resultado é uma trajetória \textbf{cicloidal}, com deriva média
\[
v_d=\frac{E}{B}=\frac{500}{0{,}02}=\pot{2.5}{4}\,\si{\meter\per\second},
\]
perpendicular a ambos os campos.\\
\textbf{Ponto essencial:} a força magnética \emph{depende da velocidade} --- ela
não age sobre cargas paradas. Já a elétrica age sempre. É essa assimetria que
faz o movimento começar de um jeito e evoluir para outro.}
\end{questao}

\begin{questao}[2]{}
Um fio em forma de semicircunferência de raio $R$ conduz corrente $I$ em um
campo uniforme $\vec B$ perpendicular ao plano do semicírculo. Determine a
força resultante.
\resposta{$F=2RIB$, como se o fio fosse retilíneo de comprimento $2R$}
\resolucao{Para um fio de forma qualquer em campo \textbf{uniforme},
\[
\vec F=I\left(\int\mathrm{d}\vec\ell\right)\times\vec B
=I\,\vec L_{ef}\times\vec B,
\]
onde $\vec L_{ef}$ é o vetor que liga o \textbf{início ao fim} do fio --- a
\emph{corda}, não o comprimento percorrido.\\
Para a semicircunferência, a corda vale $2R$:
\[
F=2RIB.
\]
\textbf{Poder do resultado:} a forma do fio é irrelevante em campo uniforme.
Um fio retorcido, ondulado ou em zigue-zague sofre a mesma força que o segmento
reto entre suas extremidades.\\
\textbf{Corolário importante:} em uma \emph{espira fechada}, a corda é nula ---
logo a força resultante sobre qualquer espira em campo uniforme é
\textbf{zero}. Só há torque.}
\end{questao}

\begin{questao}[2]{}
Uma balança de corrente tem um fio horizontal de $\SI{20}{\centi\meter}$ em um
campo vertical de $\SI{0.4}{\tesla}$. Determine a corrente necessária para
equilibrar uma massa de $\SI{2}{\gram}$.
\dados{$g=\SI{9.8}{\meter\per\second\squared}$}
\resposta{$I=\SI{0.245}{\ampere}$}
\resolucao{Equilíbrio entre peso e força magnética:
\[
BIL=mg
\;\Longrightarrow\;
I=\frac{mg}{BL}=\frac{(\pot{2}{-3})(9{,}8)}{(0{,}4)(0{,}20)}
=\SI{0.245}{\ampere}.
\]
\textbf{Uso do arranjo:} conhecendo $m$, $L$ e $I$, obtém-se $B$ com boa
exatidão --- é o método clássico de calibração absoluta de campo magnético, sem
depender de sensores.\\
\textbf{Inversamente}, conhecendo $B$, mede-se corrente. Foi assim que o ampère
foi definido durante décadas: pela força entre condutores.}
\end{questao}

\begin{questao}[2]{}
Determine a velocidade de uma partícula que atravessa sem desvio uma região com
$E=\SI{2e4}{\volt\per\meter}$ e $B=\SI{0.1}{\tesla}$ perpendiculares entre si
e à velocidade.
\resposta{$v=\pot{2}{5}\,\si{\meter\per\second}$}
\resolucao{Para não haver desvio, as forças elétrica e magnética devem se
cancelar:
\[
qE=qvB
\;\Longrightarrow\;
v=\frac{E}{B}=\frac{\pot{2}{4}}{0{,}1}=\pot{2}{5}\,\si{\meter\per\second}.
\]
\textbf{Duas propriedades notáveis:}
\begin{itemize}
\item A condição \textbf{não depende da carga} --- $q$ cancela. Íons e elétrons
com a mesma velocidade passam igualmente.
\item Também \textbf{não depende da massa}. É pura seleção de \emph{velocidade}.
\end{itemize}
\textbf{Consequência:} partículas mais rápidas são desviadas pela força
magnética (maior), e mais lentas pela elétrica. Apenas a velocidade $E/B$
atravessa em linha reta.}
\end{questao}

\begin{questao}[2]{}
Uma espira quadrada de lado $\SI{10}{\centi\meter}$ conduz $\SI{5}{\ampere}$
em um campo uniforme de $\SI{0.2}{\tesla}$ paralelo ao seu plano.
\begin{itensq}
\item Determine a força resultante.
\item Determine o torque.
\end{itensq}
\resposta{$\vec F=\vec 0$; $\tau=\SI{10}{\milli\newton\meter}$}
\resolucao{(a) \textbf{Força resultante nula.} Em campo uniforme, os quatro
lados sofrem forças que se cancelam duas a duas --- lados opostos conduzem em
sentidos contrários.\\
(b) \textbf{Mas o torque não é nulo.} As forças nos dois lados perpendiculares
ao campo formam um \emph{binário}:
\[
\tau=mB\sin\ang{90}=IAB=(5)(\pot{1}{-2})(0{,}2)
=\pot{1}{-2}\,\si{\newton\meter}.
\]
\textbf{Situação de torque máximo:} o campo está no plano da espira, ou seja,
$\vec m$ (perpendicular ao plano) está a $\ang{90}$ de $\vec B$.\\
\textbf{Este é o princípio do motor elétrico} --- e também sua limitação: o
torque cai a zero quando a espira gira até o alinhamento, o que exige comutador
para inverter a corrente e manter a rotação.}
\end{questao}

\blocoaprofundamento

\begin{questao}[3]{}
Uma partícula de carga $q$ e massa $m$ entra perpendicularmente em um campo
uniforme $B$ com velocidade $v$.
\begin{itensq}
\item Deduza o raio e o período da trajetória.
\item Determine a frequência angular (frequência de cíclotron).
\item Explique por que o período independe de $v$.
\end{itensq}
\resposta{$r=\dfrac{mv}{qB}$; $T=\dfrac{2\pi m}{qB}$;
$\omega=\dfrac{qB}{m}$}
\resolucao{(a) A força magnética é perpendicular a $\vec v$ e de módulo
constante --- exatamente as condições de um movimento \textbf{circular
uniforme}. Igualando à resultante centrípeta:
\[
qvB=\frac{mv^{2}}{r}
\;\Longrightarrow\;
r=\frac{mv}{qB}.
\]
\[
T=\frac{2\pi r}{v}=\frac{2\pi m}{qB}.
\]
(b)
\[
\omega=\frac{2\pi}{T}=\frac{qB}{m},
\]
a \textbf{frequência angular de cíclotron}.\\
(c) \textbf{Por que $T$ independe de $v$.} Dobrar a velocidade dobra a força
magnética \emph{e} dobra o raio necessário --- os dois efeitos se compensam
exatamente. A partícula percorre um círculo duas vezes maior a velocidade duas
vezes maior, no mesmo tempo.\\
Algebricamente, $v$ aparece em $r$ e é cancelado por $v$ em $T=2\pi r/v$.\\
\textbf{Consequência:} o cíclotron pode operar com frequência \emph{fixa}.
\emph{(Isso vale enquanto $v\ll c$; em regime relativístico a massa efetiva
cresce, $T$ aumenta e é preciso modular a frequência --- o síncrotron.)}}
\end{questao}

\begin{questao}[3]{}
Uma partícula entra em um campo uniforme formando $\ang{30}$ com $\vec B$, com
$v=\pot{1}{6}\,\si{\meter\per\second}$ e $B=\SI{0.1}{\tesla}$ (próton).
\begin{itensq}
\item Decomponha a velocidade e descreva a trajetória.
\item Determine o raio e o passo da hélice.
\end{itensq}
\resposta{Hélice com $r=\SI{5.2}{\centi\meter}$ e passo de
$\SI{0.57}{\meter}$}
\resolucao{(a) \textbf{Decomposição:}
\[
v_\perp=v\sin\ang{30}=\pot{5}{5}\,\si{\meter\per\second};
\qquad
v_\parallel=v\cos\ang{30}=\pot{8.66}{5}\,\si{\meter\per\second}.
\]
A componente \textbf{perpendicular} sofre força e produz movimento circular. A
componente \textbf{paralela} não sofre força alguma ($\vec v\times\vec B=0$
para ela) e produz movimento retilíneo uniforme.\\
A composição é uma \textbf{hélice} de eixo paralelo a $\vec B$.\\
(b)
\[
r=\frac{mv_\perp}{qB}
=\frac{(\pot{1.67}{-27})(\pot{5}{5})}{(\pot{1.6}{-19})(0{,}1)}
=\SI{0.052}{\meter};
\]
\[
T=\frac{2\pi m}{qB}=\pot{6.56}{-7}\,\si{\second};
\qquad
\text{passo}=v_\parallel T=(\pot{8.66}{5})(\pot{6.56}{-7})=\SI{0.568}{\meter}.
\]
\textbf{Aplicação:} é essa trajetória que confina partículas nos cinturões de
Van Allen e em espelhos magnéticos --- as partículas espiralam ao longo das
linhas de campo, e o passo diminui onde o campo se intensifica, até a reversão
do movimento.}
\end{questao}

\begin{questao}[3]{}
Um espectrômetro de massa usa um seletor de velocidades ($v=\pot{1}{5}\,
\si{\meter\per\second}$) seguido de um campo de $\SI{0.1}{\tesla}$. Íons de
neônio de massas $20\,u$ e $22\,u$, todos com carga $+e$, entram no segundo
campo.
\begin{itensq}
\item Determine o raio de cada trajetória.
\item Determine a separação entre os pontos de chegada no detector.
\end{itensq}
\dados{$u=\pot{1.66}{-27}\,\si{\kilogram}$}
\resposta{$r_{20}=\SI{20.8}{\centi\meter}$, $r_{22}=\SI{22.8}{\centi\meter}$;
separação de $\SI{4.2}{\centi\meter}$}
\resolucao{(a) Com $r=mv/(qB)$ e velocidade já selecionada:
\[
r_{20}=\frac{20(\pot{1.66}{-27})(\pot{1}{5})}{(\pot{1.6}{-19})(0{,}1)}
=\SI{0.2075}{\meter};
\qquad
r_{22}=\SI{0.2283}{\meter}.
\]
(b) Após meia volta, cada íon atinge o detector a uma distância $2r$ do ponto
de entrada:
\[
\Delta=2\left(r_{22}-r_{20}\right)=2(0{,}0208)=\SI{4.15}{\centi\meter}.
\]
\textbf{Por que o seletor é indispensável:} sem ele, íons de massas diferentes
mas velocidades diferentes poderiam produzir o \emph{mesmo} raio --- pois
$r\propto mv$. O seletor fixa $v$, tornando $r$ função apenas de $m/q$.\\
\textbf{Resolução do instrumento:} a separação relativa é
$\Delta r/r=\Delta m/m=10\%$ para esses isótopos --- larga. Separar isótopos
com $\Delta m/m$ de $0{,}1\%$ exigiria feixe muito bem colimado e detector de
alta resolução espacial.}
\end{questao}

\begin{questao}[3]{}
A definição do ampère baseava-se na força entre condutores paralelos.
\begin{itensq}
\item Determine a força por metro entre dois fios distantes
      $\SI{1}{\meter}$, conduzindo $\SI{1}{\ampere}$ cada.
\item Explique o papel dessa relação na definição da unidade.
\item Comente a mudança recente.
\end{itensq}
\resposta{$F/L=\pot{2}{-7}\,\si{\newton\per\meter}$}
\resolucao{(a)
\[
\frac{F}{L}=\frac{\mu_0I^{2}}{2\pi d}
=\frac{(4\pi\times10^{-7})(1)}{2\pi(1)}=\pot{2}{-7}\,\si{\newton\per\meter}.
\]
(b) \textbf{Definição antiga (até 2019):} o ampère era \emph{definido} como a
corrente que, em dois condutores paralelos infinitos separados por
$\SI{1}{\meter}$ no vácuo, produzisse exatamente
$\pot{2}{-7}\,\si{\newton}$ por metro.\\
Essa definição \textbf{fixava $\mu_0$} em $4\pi\times10^{-7}$ por construção ---
não era um valor medido, mas uma consequência da escolha.\\
(c) \textbf{Desde 2019}, o ampère é definido fixando-se a \emph{carga
elementar} em $e=\pot{1.602176634}{-19}\,\si{\coulomb}$ exatamente.\\
\textbf{Consequência da mudança:} $\mu_0$ deixou de ser exato e passou a ser uma
grandeza \textbf{medida}, com incerteza experimental --- embora ainda
igual a $4\pi\times10^{-7}$ dentro de partes por bilhão.\\
A troca reflete uma mudança de filosofia: definir unidades por
\emph{constantes fundamentais} em vez de por experimentos difíceis de
reproduzir.}
\end{questao}

\begin{questao}[3]{}
Um sensor Hall de espessura $\SI{0.1}{\milli\meter}$ conduz
$\SI{1}{\ampere}$ em campo de $\SI{0.5}{\tesla}$.
\begin{itensq}
\item Determine a tensão Hall para cobre ($n=\pot{8.5}{28}\,\si{\per\meter\cubed}$).
\item Repita para um semicondutor com $n=\pot{1}{22}\,\si{\per\meter\cubed}$.
\item Explique a diferença.
\end{itensq}
\resposta{$\SI{0.37}{\micro\volt}$ no cobre; $\SI{3.1}{\volt}$ no
semicondutor}
\resolucao{(a) A tensão Hall vale
\[
V_H=\frac{IB}{nqt}
=\frac{(1)(0{,}5)}{(\pot{8.5}{28})(\pot{1.6}{-19})(\pot{1}{-4})}
=\pot{3.68}{-7}\,\si{\volt}.
\]
(b) Com $n=\pot{1}{22}$:
\[
V_H=\frac{0{,}5}{(\pot{1}{22})(\pot{1.6}{-19})(\pot{1}{-4})}=\SI{3.1}{\volt}.
\]
\textbf{Diferença de sete ordens de grandeza.}\\
(c) \textbf{Explicação.} A tensão Hall é inversamente proporcional a $n$. Metais
têm um portador por átomo --- densidade altíssima --- e cada portador precisa se
mover \emph{devagar} para produzir a mesma corrente. Velocidade de deriva baixa
significa força magnética pequena e, portanto, separação de cargas mínima.\\
Semicondutores têm densidade de portadores milhões de vezes menor; eles se movem
muito mais rápido, e o efeito Hall se torna macroscópico.\\
\textbf{Consequência tecnológica:} \emph{todos} os sensores Hall comerciais são
de semicondutor. Um sensor de cobre exigiria amplificação de nanovolts,
inviável na presença de ruído térmico e de tensões termoelétricas.}
\end{questao}

\begin{questao}[3]{}
Um motor elementar tem espira de $200$ voltas, área
$\SI{50}{\centi\meter\squared}$, corrente de $\SI{3}{\ampere}$ e campo de
$\SI{0.4}{\tesla}$.
\begin{itensq}
\item Determine o torque máximo.
\item Determine a potência mecânica a $\SI{1200}{rpm}$, supondo torque médio
      igual a $2/\pi$ do máximo.
\end{itensq}
\resposta{$\tau_{\max}=\SI{1.2}{\newton\meter}$;
$P\approx\SI{96}{\watt}$}
\resolucao{(a)
\[
m=NIA=200(3)(\pot{5}{-3})=\SI{3}{\ampere\meter\squared};
\qquad
\tau_{\max}=mB=3(0{,}4)=\SI{1.2}{\newton\meter}.
\]
\emph{(Atenção à conversão: $\SI{50}{\centi\meter\squared}
=\pot{5}{-3}\,\si{\meter\squared}$, e não $\pot{5}{-2}$.)}\\
(b) O torque varia como $\sin\theta$ ao longo da volta, e seu valor médio sobre
meio ciclo é $\dfrac{2}{\pi}\tau_{\max}=\SI{0.764}{\newton\meter}$.
\[
\omega=1200\ \text{rpm}=\frac{1200(2\pi)}{60}=\SI{125.7}{\radian\per\second};
\]
\[
P=\tau_{\text{méd}}\,\omega=0{,}764(125{,}7)=\SI{96}{\watt}.
\]
\textbf{Por que o comutador é necessário:} sem inverter a corrente a cada meia
volta, o torque mudaria de sinal e a espira apenas oscilaria em torno da posição
de alinhamento.\\
\textbf{Por que motores reais usam muitas bobinas:} com uma só, o torque
\emph{pulsa} entre zero e o máximo. Várias bobinas defasadas produzem torque
quase constante --- e é isso que se busca em qualquer máquina prática.}
\end{questao}

\begin{questao}[3]{}
Uma espira de corrente é colocada em um campo magnético \textbf{não uniforme},
mais intenso de um lado.
\begin{itensq}
\item Determine se há força resultante.
\item Determine o sentido dessa força.
\item Relacione com o comportamento de ímãs.
\end{itensq}
\resposta{Há força resultante, dirigida para a região de campo mais intenso (se
$\vec m$ estiver alinhado com $\vec B$)}
\resolucao{(a) \textbf{Sim.} Em campo \emph{uniforme}, as forças nos lados
opostos da espira se cancelam exatamente e resta apenas torque. Em campo
\emph{não uniforme}, os lados estão em campos de módulos diferentes --- e o
cancelamento é incompleto.\\
(b) \textbf{Sentido.} A força resultante é dada pelo gradiente:
\[
\vec F=\nabla\left(\vec m\cdot\vec B\right).
\]
Se $\vec m$ estiver \emph{alinhado} com $\vec B$, a espira é atraída para a
região de campo mais intenso. Se estiver \emph{antialinhado}, é repelida.\\
(c) \textbf{Relação com ímãs.} Um ímã permanente é, essencialmente, um conjunto
de espiras microscópicas (correntes de magnetização) com momentos alinhados.
Daí:
\begin{itemize}
\item \textbf{Dois ímãs se atraem ou repelem} conforme a orientação relativa dos
momentos --- exatamente o comportamento da espira.
\item \textbf{Um ímã atrai ferro não magnetizado} porque primeiro o magnetiza
(alinhando seus domínios com o campo aplicado) e depois o puxa para a região de
campo mais forte. A atração é sempre no sentido do gradiente.
\item \textbf{Em campo uniforme, um ímã não é atraído} --- ele apenas gira até
se alinhar. É por isso que a agulha de uma bússola aponta, mas não é puxada em
direção ao norte.
\end{itemize}}
\end{questao}

\begin{questao}[3]{}
Compare a força entre dois fios paralelos conduzindo correntes iguais no mesmo
sentido e em sentidos opostos.
\begin{itensq}
\item Determine o sentido em cada caso.
\item Determine o módulo.
\item Explique a diferença em termos de campo.
\end{itensq}
\resposta{Mesmo sentido: atração; opostos: repulsão. Módulos iguais}
\resolucao{(a) e (b) O \textbf{módulo} é o mesmo nos dois casos:
\[
\frac{F}{L}=\frac{\mu_0I_1I_2}{2\pi d}.
\]
O que muda é o \emph{sentido}: correntes paralelas \textbf{atraem-se},
antiparalelas \textbf{repelem-se}.\\
(c) \textbf{Explicação pelo campo.} Considere o fio 2 imerso no campo do fio 1.
\begin{itemize}
\item \textbf{Mesmo sentido:} na posição do fio 2, o campo de 1 aponta em certo
sentido; aplicando $\vec F=I\vec\ell\times\vec B$, a força aponta \emph{para} o
fio 1.
\item \textbf{Sentidos opostos:} $\vec\ell$ inverte, e a força inverte junto.
\end{itemize}
\textbf{Leitura pelas linhas de campo.} Com correntes paralelas, o campo se
\emph{cancela} entre os fios e se \emph{reforça} fora --- as linhas envolvem o
par, e a "tensão" ao longo delas puxa os fios um contra o outro. Com correntes
opostas, ocorre o inverso: campo intenso entre os fios, que os afasta.\\
\textbf{Consequência prática:} em curtos-circuitos, barramentos paralelos
sofrem forças enormes. Com $\SI{20}{\kilo\ampere}$ a $\SI{10}{\centi\meter}$,
a força chega a $\SI{800}{\newton\per\meter}$ --- razão de existirem
espaçadores mecânicos robustos em painéis de potência.}
\end{questao}

\begin{questao}[3]{}
Um próton atravessa uma região com $E=\SI{2e4}{\volt\per\meter}$ e
$B=\SI{0.1}{\tesla}$, perpendiculares entre si e à velocidade, dispostos de
modo que as forças se oponham.
\begin{itensq}
\item Determine as duas parcelas da força de Lorentz para
      $v=\pot{1}{5}$, $\pot{2}{5}$ e $\pot{4}{5}\,\si{\meter\per\second}$.
\item Identifique qual parcela domina em cada caso.
\item Interprete o resultado.
\end{itensq}
\resposta{Domina a elétrica em $\pot{1}{5}$, equilíbrio em $\pot{2}{5}$ e a
magnética em $\pot{4}{5}\,\si{\meter\per\second}$}
\resolucao{(a) A parcela elétrica \textbf{não depende de $v$}:
\[
F_E=qE=(\pot{1.6}{-19})(\pot{2}{4})=\pot{3.2}{-15}\,\si{\newton}.
\]
A magnética é \textbf{proporcional a $v$}: $F_B=qvB$.
\begin{center}
\begin{tabular}{cccl}
\hline
$v$ ($\si{\meter\per\second}$) & $F_E$ ($\si{\newton}$)
 & $F_B$ ($\si{\newton}$) & Resultante\\
\hline
$\pot{1}{5}$ & $\pot{3.2}{-15}$ & $\pot{1.6}{-15}$ & elétrica domina\\
$\pot{2}{5}$ & $\pot{3.2}{-15}$ & $\pot{3.2}{-15}$ & \textbf{nula}\\
$\pot{4}{5}$ & $\pot{3.2}{-15}$ & $\pot{6.4}{-15}$ & magnética domina\\
\hline
\end{tabular}
\end{center}
(b) e (c) \textbf{Interpretação.} O equilíbrio ocorre em
\[
v=\frac{E}{B}=\frac{\pot{2}{4}}{0{,}1}=\pot{2}{5}\,\si{\meter\per\second},
\]
e apenas nessa velocidade a partícula segue reta.\\
\textbf{As três situações da tabela descrevem o mesmo dispositivo}: partículas
lentas são empurradas para um lado pela força elétrica, rápidas para o outro
pela magnética, e só as de velocidade $E/B$ atravessam sem desvio. É a
\emph{seleção de velocidades}, vista agora como competição entre as duas
parcelas de $\vec F=q(\vec E+\vec v\times\vec B)$.\\
\textbf{Observe a estrutura:} uma parcela constante e outra linear em $v$. A
igualdade só pode ocorrer em \emph{um} valor de velocidade --- e é isso que dá
seletividade ao instrumento.}
\end{questao}

\begin{questao}[3]{}
Uma partícula carregada move-se em uma região com campos $\vec E$ e $\vec B$
perpendiculares entre si, mas com velocidade \emph{diferente} de $E/B$.
\begin{itensq}
\item Descreva qualitativamente o movimento.
\item Determine a velocidade média de deriva.
\item Indique uma aplicação.
\end{itensq}
\resposta{Movimento cicloidal, com deriva média
$\vec v_d=\dfrac{\vec E\times\vec B}{B^{2}}$}
\resolucao{(a) \textbf{Movimento.} A partícula descreve uma trajetória
\textbf{cicloidal}: um movimento circular (devido a $\vec B$) superposto a um
deslocamento uniforme. Se ela parte do repouso, a curva é uma cicloide clássica
--- como um ponto na borda de uma roda que rola.\\
(b) \textbf{Velocidade de deriva.} O centro do círculo desloca-se com
\[
\vec v_d=\frac{\vec E\times\vec B}{B^{2}},
\qquad
|v_d|=\frac{E}{B},
\]
perpendicular a \emph{ambos} os campos.\\
\textbf{Duas propriedades notáveis:} a deriva \textbf{não depende da carga}
(íons e elétrons derivam no mesmo sentido) nem da \textbf{massa}.\\
(c) \textbf{Aplicações.}
\begin{itemize}
\item \textbf{Magnetrons} (fornos de micro-ondas): elétrons em campos cruzados
descrevem trajetórias cicloidais que excitam cavidades ressonantes.
\item \textbf{Confinamento de plasma:} a deriva $\vec E\times\vec B$ é um dos
mecanismos que governam o transporte em tokamaks --- e como ela não separa
cargas, não gera corrente própria.
\item \textbf{Magnetosfera terrestre:} explica parte do movimento de partículas
aprisionadas.
\end{itemize}}
\end{questao}

\begin{questao}[3]{}
Um condutor de $\SI{50}{\centi\meter}$ e $\SI{100}{\gram}$ está apoiado sobre
dois trilhos horizontais em campo vertical de $\SI{0.8}{\tesla}$. O
coeficiente de atrito estático é $0{,}25$.
\begin{itensq}
\item Determine a corrente mínima para iniciar o movimento.
\item Determine a aceleração inicial se a corrente for o dobro, com atrito
      cinético de $0{,}20$.
\end{itensq}
\resposta{$I=\SI{0.613}{\ampere}$; $a\approx\SI{2.9}{\meter\per\second\squared}$}
\resolucao{(a) A força magnética deve vencer o atrito estático:
\[
BIL=\mu_e mg
\;\Longrightarrow\;
I=\frac{0{,}25(0{,}100)(9{,}8)}{(0{,}8)(0{,}50)}=\SI{0.613}{\ampere}.
\]
(b) Com $I=\SI{1.226}{\ampere}$:
\[
F=BIL=(0{,}8)(1{,}226)(0{,}50)=\SI{0.490}{\newton};
\]
\[
f_c=\mu_c mg=0{,}20(0{,}100)(9{,}8)=\SI{0.196}{\newton};
\]
\[
a=\frac{0{,}490-0{,}196}{0{,}100}=\SI{2.94}{\meter\per\second\squared}.
\]
\textbf{Este é o princípio do canhão eletromagnético (\emph{railgun})} --- e
também do motor linear. \\
\textbf{Ressalva importante:} conforme o condutor acelera, ele passa a induzir
f.e.m. contrária ($\varepsilon=BLv$), que reduz a corrente. A aceleração
calculada vale apenas no \emph{instante inicial}; o movimento tende a uma
velocidade-limite em que a f.e.m. induzida equilibra a fonte.}
\end{questao}

\blocodesafio

\begin{questao}[4]{}
\textbf{(Demonstração --- trabalho nulo.)} Prove que a força magnética nunca
realiza trabalho sobre uma carga puntiforme.
\begin{itensq}
\item Demonstre formalmente.
\item Concilie com o fato de que motores elétricos realizam trabalho.
\item Explique o que ocorre com a energia em um motor.
\end{itensq}
\resposta{$\vec F\cdot\vec v=0$ sempre; o trabalho no motor vem da fonte, não do
campo}
\resolucao{(a) A potência é
\[
P=\vec F\cdot\vec v=q\left(\vec v\times\vec B\right)\cdot\vec v.
\]
O produto misto $(\vec v\times\vec B)\cdot\vec v$ é \textbf{identicamente
nulo}, pois $\vec v\times\vec B$ é perpendicular a $\vec v$ por construção. Logo
$P=0$ e $\Delta E_c=0$: o \emph{módulo} da velocidade nunca muda.\\
(b) \textbf{O aparente paradoxo do motor.} Um motor claramente realiza trabalho
--- ele levanta cargas e move veículos. Como conciliar?\\
\textbf{Resolução:} a força magnética atua sobre os \emph{portadores} do
condutor, e é perpendicular à velocidade \emph{deles}. Mas o fio se move, e a
velocidade total de um portador é a soma da deriva ao longo do fio com a
velocidade do fio.\\
A força magnética \emph{redistribui} energia: ela empurra o fio (realizando
trabalho positivo sobre ele) e, simultaneamente, freia a deriva dos portadores
(trabalho negativo de mesmo módulo). A soma continua sendo zero.\\
(c) \textbf{De onde vem a energia.} O freio à deriva manifesta-se como
\textbf{f.e.m. contraeletromotriz}, e é a \emph{fonte} que precisa vencê-la:
\[
P_{fonte}=\varepsilon_{fem}\,I=P_{\text{mecânica}}+P_{Joule}.
\]
\textbf{O campo magnético é um intermediário, não uma fonte.} Ele transfere
energia da fonte elétrica para o eixo, sem fornecer nem armazenar nada
líquido.\\
\emph{(É o mesmo papel de uma polia em mecânica: ela redireciona força sem criar
energia.)}}
\end{questao}

\begin{questao}[4]{}
\textbf{(Dedução do torque.)} Deduza a expressão do torque sobre uma espira
retangular em campo uniforme.
\begin{itensq}
\item Analise as forças em cada lado.
\item Obtenha $\tau=NIAB\sin\theta$.
\item Generalize para espira de forma qualquer.
\end{itensq}
\resposta{$\vec\tau=\vec m\times\vec B$, com $\vec m=NI\vec A$}
\resolucao{(a) Seja a espira de lados $a$ e $b$, com $\vec B$ no plano
horizontal e o eixo de rotação vertical. Chamando $\theta$ o ângulo entre
$\vec B$ e a \emph{normal} à espira:
\begin{itemize}
\item Os lados de comprimento $a$ \emph{paralelos} ao eixo sofrem forças
$F=BIa$, perpendiculares tanto ao fio quanto ao campo, em sentidos opostos.
Elas formam um \textbf{binário}.
\item Os lados de comprimento $b$ sofrem forças que se cancelam e cuja linha de
ação passa pelo eixo --- não produzem torque.
\end{itemize}
(b) O braço de alavanca de cada força do binário é
$\dfrac{b}{2}\sin\theta$:
\[
\tau=2\left(BIa\right)\frac{b}{2}\sin\theta=BI(ab)\sin\theta=BIA\sin\theta.
\]
Com $N$ espiras, $\tau=NIAB\sin\theta$.\\
(c) \textbf{Generalização.} Definindo o momento magnético
$\vec m=NI\vec A$ (com $\vec A$ normal à espira, sentido dado pela regra da mão
direita):
\[
\boxed{\vec\tau=\vec m\times\vec B}
\]
Essa forma vale para espira de \textbf{qualquer} formato --- basta decompô-la em
retângulos infinitesimais, cujas correntes internas se cancelam aos pares.\\
\textbf{Analogia formal com o dipolo elétrico:} $\vec\tau=\vec p\times\vec E$.
As duas expressões são idênticas em estrutura, e por isso a espira é
literalmente um \emph{dipolo magnético}.}
\end{questao}

\begin{questao}[4]{}
\textbf{(Projeto de espectrômetro.)} Projete um espectrômetro capaz de separar
isótopos com $\Delta m/m=1\%$.
\begin{itensq}
\item Determine a separação espacial em função dos parâmetros.
\item Estabeleça o requisito de colimação do feixe.
\item Discuta as limitações práticas.
\end{itensq}
\resposta{$\Delta r/r=\Delta m/m$; a resolução exige feixe com dispersão angular
e de velocidade menores que $1\%$}
\resolucao{(a) De $r=\dfrac{mv}{qB}$ com $v$ e $B$ fixos:
\[
\frac{\Delta r}{r}=\frac{\Delta m}{m}=1\%.
\]
Com $r=\SI{20}{\centi\meter}$, a separação de raios é
$\SI{2}{\milli\meter}$, e no detector (após meia volta) a separação é
$2\Delta r=\SI{4}{\milli\meter}$.\\
(b) \textbf{Requisito de colimação.} Para \emph{resolver} duas linhas separadas
por $\SI{4}{\milli\meter}$, a largura de cada linha deve ser bem menor. Duas
fontes de alargamento:
\begin{itemize}
\item \textbf{Dispersão de velocidade:} como $r\propto v$, uma dispersão
$\Delta v/v$ produz alargamento $\Delta r/r$ de mesma magnitude. É preciso
$\Delta v/v\ll1\%$ --- daí a necessidade do seletor de velocidades.
\item \textbf{Dispersão angular:} íons entrando com pequenos ângulos diferentes
descrevem círculos deslocados. \emph{Felizmente}, a geometria de meia volta é
\textbf{autofocalizante}: íons que divergem de um ponto reconvergem após
$\ang{180}$, em primeira ordem. O alargamento residual é de segunda ordem no
ângulo.
\end{itemize}
(c) \textbf{Limitações práticas.}
\begin{itemize}
\item \textbf{Vácuo:} colisões com gás residual espalham o feixe. É preciso
pressão abaixo de $\pot{1}{-5}\,\si{\pascal}$.
\item \textbf{Uniformidade do campo:} variações de $B$ ao longo da trajetória
introduzem erro proporcional. Exige-se uniformidade melhor que a resolução
desejada.
\item \textbf{Carga espacial:} feixes intensos se repelem eletrostaticamente e
se alargam --- limitando a corrente utilizável.
\item \textbf{Íons de mesma razão $m/q$:} não são separáveis. $\mathrm{N_2^+}$ e
$\mathrm{CO^+}$ (ambos $28\,u$) exigem instrumentos de altíssima resolução para
distinguir.
\end{itemize}}
\end{questao}

\begin{questao}[4]{}
\textbf{(Projeto de atuador linear.)} Projete um atuador que produza
$\SI{5}{\newton}$ com curso de $\SI{2}{\centi\meter}$.
\begin{itensq}
\item Determine as combinações de $B$, $I$ e $L$ possíveis.
\item Avalie o compromisso entre corrente e número de espiras.
\item Determine a potência dissipada e o trabalho realizado.
\end{itensq}
\resposta{Com $B=\SI{0.5}{\tesla}$ e $L=\SI{0.2}{\meter}$: $I=\SI{50}{\ampere}$
com uma espira, ou $\SI{2.5}{\ampere}$ com vinte}
\resolucao{(a) De $F=NBIL$:
\[
NIL=\frac{F}{B}=\frac{5}{0{,}5}=\SI{10}{\ampere\meter}.
\]
Qualquer combinação que satisfaça esse produto serve --- há dois graus de
liberdade.\\
(b) \textbf{Compromisso.} Com $L=\SI{0.2}{\meter}$ fixo:
\begin{center}
\begin{tabular}{ccc}
\hline
$N$ & $I$ & Observação\\
\hline
$1$ & $\SI{50}{\ampere}$ & fio muito grosso, fonte de alta corrente\\
$20$ & $\SI{2.5}{\ampere}$ & equilíbrio razoável\\
$200$ & $\SI{0.25}{\ampere}$ & bobina volumosa, alta indutância\\
\hline
\end{tabular}
\end{center}
\textbf{Mais espiras} reduzem a corrente (fio mais fino, fonte mais simples),
mas aumentam volume, massa, resistência e \emph{indutância} --- o que torna o
atuador mais lento a responder.\\
\textbf{Menos espiras} dão resposta rápida, ao custo de corrente elevada.\\
(c) \textbf{Dissipação.} Com $N=20$, $I=\SI{2.5}{\ampere}$ e resistência total
estimada em $\SI{2}{\ohm}$:
\[
P=RI^{2}=2(6{,}25)=\SI{12.5}{\watt}.
\]
\textbf{Trabalho útil:}
\[
W=F\,d=5(0{,}02)=\SI{0.1}{\joule}.
\]
\textbf{Comparação reveladora:} se o curso levar $\SI{0.1}{\second}$, a potência
mecânica é $\SI{1}{\watt}$ contra $\SI{12.5}{\watt}$ dissipados --- rendimento
de $7\%$. Atuadores eletromagnéticos de curso curto são intrinsecamente
ineficientes, e é por isso que costumam operar em regime \emph{pulsado}, e não
contínuo.}
\end{questao}

\begin{questao}[4]{}
\textbf{(Efeito Hall.)} Deduza a expressão da tensão Hall e analise suas
aplicações.
\begin{itensq}
\item Deduza $V_H=\dfrac{IB}{nqt}$.
\item Explique como o sinal revela o tipo de portador.
\item Cite três aplicações e o que cada uma mede.
\end{itensq}
\resposta{Equilíbrio entre força magnética e força elétrica transversal}
\resolucao{(a) \textbf{Dedução.} Os portadores, com velocidade de deriva
$v_d$, sofrem força magnética $qv_dB$ transversal e acumulam-se em uma face.
Isso cria campo elétrico transversal $E_H$, que cresce até equilibrar:
\[
qE_H=qv_dB
\;\Longrightarrow\;
E_H=v_dB.
\]
Com largura $w$, a tensão é $V_H=E_Hw=v_dBw$. Usando
$I=nqv_d(wt)$, elimina-se $v_d$:
\[
v_d=\frac{I}{nqwt}
\;\Longrightarrow\;
V_H=\frac{IB}{nqt}.
\]
\textbf{Note:} a largura $w$ cancela --- só a \emph{espessura} $t$ importa.\\
(b) \textbf{Sinal e tipo de portador.} Elétrons (negativos) e lacunas
(positivas) movem-se em sentidos opostos para produzir a mesma corrente. A força
$q\vec v\times\vec B$ os desvia para o \textbf{mesmo} lado --- mas com cargas de
sinais contrários, a polaridade de $V_H$ \textbf{se inverte}.\\
\textbf{É o único método simples de determinar experimentalmente o sinal dos
portadores} --- e foi assim que se estabeleceu que semicondutores tipo P
conduzem por lacunas.\\
(c) \textbf{Três aplicações.}
\begin{itemize}
\item \textbf{Medição de campo magnético} (gaussímetro): com $I$, $n$ e
$t$ conhecidos, $V_H\propto B$.
\item \textbf{Medição de corrente sem contato:} coloca-se o sensor no entreferro
de um núcleo que envolve o condutor. Mede $B$, e daí $I$ --- com isolação
galvânica completa.
\item \textbf{Caracterização de materiais:} com $B$ e $I$ conhecidos, obtém-se
$n$ (densidade de portadores) e o \emph{sinal} deles. Combinado com a
condutividade, fornece a mobilidade.
\end{itemize}
\textbf{Vantagem sobre alternativas:} sensores Hall respondem a campo
\emph{estático}, ao contrário de bobinas, que só detectam variação.}
\end{questao}

\begin{questao}[4]{}
\textbf{(Estabilidade de espira.)} Analise o equilíbrio de uma espira de
corrente em campo magnético uniforme.
\begin{itensq}
\item Determine as posições de equilíbrio.
\item Classifique cada uma quanto à estabilidade.
\item Deduza o período de pequenas oscilações.
\end{itensq}
\resposta{$\theta=0$ é estável, $\theta=\pi$ é instável;
$T=2\pi\sqrt{J/(mB)}$}
\resolucao{(a) O torque é $\tau=-mB\sin\theta$, com $\theta$ o ângulo entre
$\vec m$ e $\vec B$. Ele se anula em
\[
\theta=0 \quad\text{e}\quad \theta=\pi.
\]
(b) \textbf{Classificação pela energia.} A energia potencial de um dipolo
magnético é
\[
U=-\vec m\cdot\vec B=-mB\cos\theta.
\]
\begin{itemize}
\item $\theta=0$ ($\vec m$ paralelo a $\vec B$): $U=-mB$, \textbf{mínimo}
$\Rightarrow$ equilíbrio \textbf{estável}. Deslocada, a espira sofre torque
restaurador.
\item $\theta=\pi$ (antiparalelo): $U=+mB$, \textbf{máximo} $\Rightarrow$
equilíbrio \textbf{instável}. Qualquer perturbação a faz girar $\ang{180}$.
\end{itemize}
(c) \textbf{Pequenas oscilações.} Para $\theta\ll1$, $\sin\theta\approx\theta$ e
\[
J\ddot\theta=-mB\theta
\;\Longrightarrow\;
\ddot\theta+\frac{mB}{J}\theta=0,
\]
equação do MHS com
\[
\omega=\sqrt{\frac{mB}{J}};
\qquad
T=2\pi\sqrt{\frac{J}{mB}},
\]
onde $J$ é o momento de inércia da espira.\\
\textbf{Aplicação prática:} é assim que funcionam o \textbf{galvanômetro de
bobina móvel} --- em que uma mola acrescenta torque restaurador proporcional a
$\theta$, tornando a deflexão proporcional à corrente --- e a \textbf{bússola},
cujo período de oscilação permite medir o campo terrestre se $J$ e $m$ forem
conhecidos.\\
\textbf{Diferença de $U$ entre os dois equilíbrios:} $2mB$ --- energia
necessária para inverter a espira, e a mesma que um ímã precisa receber para ter
sua polarização revertida.}
\end{questao}

\begin{questao}[4]{}
\textbf{(Força de Lorentz --- estrutura e consequências.)} Analise a expressão
completa $\vec F=q\left(\vec E+\vec v\times\vec B\right)$.
\begin{itensq}
\item Compare as duas parcelas quanto à dependência da velocidade, à direção e
      ao trabalho realizado.
\item Avalie um próton com $v=\pot{3}{4}\,\si{\meter\per\second}$ em
      $E=\SI{500}{\volt\per\meter}$ e $B=\SI{0.02}{\tesla}$
      \emph{perpendiculares entre si}, com $\vec v$ paralelo a $\vec E$.
\item Explique por que aceleradores usam campo elétrico para acelerar e
      magnético para guiar.
\end{itensq}
\resposta{$F_E=\SI{8e-17}{\newton}$, $F_B=\SI{9.6e-17}{\newton}$,
perpendiculares; $|F|=\SI{1.25e-16}{\newton}$ a $\ang{50.2}$ de $\vec E$}
\resolucao{(a) \textbf{Comparação das duas parcelas.}
\begin{center}
\begin{tabular}{lll}
\hline
& Parcela elétrica & Parcela magnética\\
\hline
Expressão & $q\vec E$ & $q\,\vec v\times\vec B$\\
Depende de $v$? & \textbf{não} & \textbf{sim} (linear)\\
Age em repouso? & sim & não\\
Direção & paralela a $\vec E$ & $\perp$ a $\vec v$ \emph{e} a $\vec B$\\
Realiza trabalho? & \textbf{sim} & \textbf{nunca}\\
Altera $|v|$? & sim & não\\
\hline
\end{tabular}
\end{center}
\textbf{A linha decisiva é a do trabalho.} Como
$\left(\vec v\times\vec B\right)\cdot\vec v=0$ identicamente, a parcela
magnética tem potência nula --- ela desvia sem acelerar. A elétrica, ao
contrário, é a única capaz de variar a energia cinética.\\
(b) \textbf{Avaliação numérica.}
\[
F_E=qE=(\pot{1.6}{-19})(500)=\pot{8}{-17}\,\si{\newton};
\]
\[
F_B=qvB=(\pot{1.6}{-19})(\pot{3}{4})(0{,}02)=\pot{9.6}{-17}\,\si{\newton}.
\]
Com $\vec v$ paralelo a $\vec E$ e $\vec B$ perpendicular a ambos, a força
magnética é \textbf{perpendicular} à elétrica:
\[
|\vec F|=\sqrt{F_E^{2}+F_B^{2}}=\pot{1.25}{-16}\,\si{\newton},
\qquad
\theta=\arctan\frac{9{,}6}{8{,}0}=\ang{50.2}
\]
em relação a $\vec E$.\\
\textbf{Razão entre as parcelas:}
$\dfrac{F_B}{F_E}=\dfrac{vB}{E}=1{,}20$ --- adimensional, e igual a $1$
exatamente na velocidade $E/B$ do seletor.\\
\textbf{Trabalho em um deslocamento de $\SI{1}{\centi\meter}$ ao longo de
$\vec E$:} apenas a parcela elétrica contribui,
\[
W=F_E d=(\pot{8}{-17})(0{,}01)=\pot{8}{-19}\,\si{\joule}=\SI{5}{\electronvolt},
\]
que é simplesmente $qEd$ --- ou seja, $q$ vezes a tensão percorrida.\\
(c) \textbf{Divisão de papéis em aceleradores.}
\begin{itemize}
\item \textbf{Campo elétrico acelera.} É o único que realiza trabalho, e o ganho
de energia por passagem é $qU$ --- daí as cavidades de radiofrequência.
\item \textbf{Campo magnético guia.} Ele curva a trajetória sem custo
energético, permitindo fechar o percurso em círculo e reaproveitar as mesmas
cavidades milhares de vezes.
\item \textbf{Campo magnético também focaliza}, com quadrupolos que comprimem o
feixe transversalmente --- de novo, sem alterar a energia.
\end{itemize}
\textbf{Se fosse possível acelerar com campo magnético}, aceleradores lineares
de quilômetros seriam desnecessários. A impossibilidade não é técnica: é
consequência direta de $\left(\vec v\times\vec B\right)\perp\vec v$.\\
\textbf{Nota histórica:} a expressão completa reúne dois fenômenos que foram
descobertos e estudados separadamente. Sua unificação --- e a constatação de que
o que é campo elétrico em um referencial pode ser magnético em outro --- é um
dos pontos de partida da relatividade restrita.}
\end{questao}

\begin{questao}[4]{}
\textbf{(Projeto de cíclotron.)} Um cíclotron acelera prótons até
$\SI{10}{\mega\electronvolt}$ usando campo de $\SI{1.5}{\tesla}$.
\begin{itensq}
\item Determine a frequência de operação.
\item Determine o raio máximo e o número de voltas, com
      $\SI{50}{\kilo\volt}$ de ganho por passagem.
\item Determine se a correção relativística é necessária.
\end{itensq}
\dados{$e=\pot{1.6}{-19}\,\si{\coulomb}$; $m_p=\pot{1.67}{-27}\,\si{\kilogram}$;
$\SI{1}{\mega\electronvolt}=\pot{1.6}{-13}\,\si{\joule}$}
\resposta{$f=\SI{22.9}{\mega\hertz}$; $r=\SI{30.5}{\centi\meter}$;
$100$ voltas; correção de $1{,}1\%$}
\resolucao{(a) A frequência de cíclotron independe da energia:
\[
f=\frac{qB}{2\pi m}=\frac{(\pot{1.6}{-19})(1{,}5)}{2\pi(\pot{1.67}{-27})}
=\SI{22.9}{\mega\hertz}.
\]
\textbf{É isso que torna o aparelho viável:} o oscilador que alimenta os "dês"
opera em frequência \emph{fixa}, embora a partícula ganhe energia
continuamente.\\
(b) Da energia final,
\[
v=\sqrt{\frac{2E}{m}}
=\sqrt{\frac{2(10)(\pot{1.6}{-13})}{\pot{1.67}{-27}}}
=\pot{4.38}{7}\,\si{\meter\per\second};
\]
\[
r=\frac{mv}{qB}=\SI{0.305}{\meter}
\;\Longrightarrow\;
\text{diâmetro}\approx\SI{61}{\centi\meter}.
\]
Com dois cruzamentos do intervalo por volta, cada volta rende
$\SI{100}{\kilo\electronvolt}$:
\[
N=\frac{10\,\si{\mega\electronvolt}}{0{,}1\,\si{\mega\electronvolt}}
=100\ \text{voltas}.
\]
\textbf{Tempo total:} $100/f=\SI{4.4}{\micro\second}$.\\
(c) \textbf{Correção relativística.} Com
$v/c=0{,}146$:
\[
\gamma=\frac{1}{\sqrt{1-(v/c)^{2}}}=1{,}0108,
\]
ou seja, a massa efetiva é $1{,}1\%$ maior no fim da aceleração.\\
\textbf{Consequência:} a frequência de cíclotron cai na mesma proporção, e a
partícula começa a chegar \emph{atrasada} ao intervalo acelerador. Após muitas
voltas o defasamento acumula e a aceleração cessa.\\
\textbf{Aqui, $1{,}1\%$ é tolerável} --- o cíclotron clássico funciona. Mas para
$\SI{100}{\mega\electronvolt}$ o desvio passaria de $10\%$ e seria proibitivo.\\
\textbf{Soluções:} o \textbf{sincrociclotron} modula a frequência ao longo do
ciclo; o \textbf{síncrotron} mantém o raio fixo e varia o campo. Ambos abandonam
a simplicidade do cíclotron justamente porque o período deixa de ser constante
--- e a raiz disso está na expressão $T=2\pi m/qB$, válida apenas enquanto
$m$ o for.}
\end{questao}

\begin{questao}[4]{}
\textbf{(Força entre correntes e o limite de barramentos.)} Um barramento de
subestação tem dois condutores paralelos separados por $\SI{15}{\centi\meter}$,
que devem suportar um curto-circuito de $\SI{40}{\kilo\ampere}$.
\begin{itensq}
\item Determine a força por metro durante o curto.
\item Determine o espaçamento máximo entre suportes, para deflexão admissível.
\item Discuta o caráter dessa solicitação.
\end{itensq}
\resposta{$F/L\approx\SI{2.1}{\kilo\newton\per\meter}$ --- comparável ao peso
de uma pessoa por metro}
\resolucao{(a)
\[
\frac{F}{L}=\frac{\mu_0I^{2}}{2\pi d}
=\pot{2}{-7}\frac{(\pot{4}{4})^{2}}{0{,}15}
=\pot{2}{-7}\frac{\pot{1.6}{9}}{0{,}15}=\SI{2133}{\newton\per\meter}.
\]
Mais de duas toneladas-força por metro de barramento --- e \textbf{repulsiva},
se as correntes forem opostas (que é o caso em um curto entre fases).\\
(b) \textbf{Espaçamento entre suportes.} Tratando o barramento como viga
biapoiada com carga distribuída $w=F/L$, a flecha é
\[
\delta=\frac{5wL_s^{4}}{384EJ}.
\]
Para deflexão admissível $\delta$, o vão máximo cresce apenas com
$L_s\propto\delta^{1/4}$ --- \textbf{dependência muito fraca}. Dobrar a flecha
admissível permite aumentar o vão em apenas $19\%$.\\
Na prática, os suportes ficam a cada $\SI{0.5}{}$ a $\SI{1}{\meter}$.\\
(c) \textbf{Caráter da solicitação.}
\begin{itemize}
\item \textbf{É transitória} --- dura o tempo de atuação da proteção, tipicamente
$\SI{100}{\milli\second}$. Mas a força \emph{de pico} pode ser o dobro da
calculada, devido à componente assimétrica da corrente de curto.
\item \textbf{Vai com $I^{2}$} --- dobrar a corrente de curto quadruplica a
força. É por isso que a corrente de curto-circuito presumida é um dado de
projeto tão crítico.
\item \textbf{É dinâmica:} se a frequência natural da estrutura estiver próxima
de $\SI{120}{\hertz}$ (dobro da rede, pois $F\propto i^{2}$), pode haver
ressonância mecânica.
\end{itemize}
\textbf{Conclusão:} em regime normal ($\SI{1}{\kilo\ampere}$, digamos), a força
seria $\SI{1.3}{\newton\per\meter}$ --- desprezível. O dimensionamento mecânico
de barramentos é governado inteiramente pela condição de \emph{falta}.}
\end{questao}
